\chapter{Biology Primer for the Demo Questions}
\label{app:biology}

Interviewers sometimes try the demo and ask what an answer means. This appendix explains the biology behind the demo's examples and the gold-set questions discussed in this report. \textbf{All of it is general knowledge, not taken from the repository or from the textbooks' indexed text}, so every section is marked as background. For anything you plan to say with authority, check the corresponding section of the OpenStax books.

\begin{background}[cells and membranes]
Cells are bounded by a plasma membrane made of a phospholipid bilayer with embedded proteins. The bilayer is selectively permeable: small non-polar molecules cross it easily, while ions and large polar molecules need transport proteins.
\end{background}

\begin{background}[facilitated diffusion of glucose]
Diffusion moves substances down their concentration gradient without energy input. Glucose is too large and polar to cross the lipid bilayer directly, so it enters most cells by \emph{facilitated diffusion} through carrier proteins of the GLUT family, which bind glucose and change shape to release it on the other side. It is still passive: no ATP is used.
\end{background}

\begin{background}[mitosis and anaphase]
Mitosis divides a duplicated nucleus into two identical nuclei in stages usually named prophase, prometaphase, metaphase, anaphase and telophase. In \emph{anaphase}, the cohesin proteins holding sister chromatids together are broken down, the chromatids separate at the centromere and become individual chromosomes, and kinetochore microtubules shorten to pull them to opposite poles, while polar microtubules push the poles apart and the cell elongates.
\end{background}

\begin{background}[tissues]
A tissue is a group of similar cells, with their surrounding material, that work together on a common function. Human tissues are classified into four types: epithelial (covering and lining), connective (support and binding, including bone and blood), muscle (contraction) and nervous (signalling).
\end{background}

\begin{background}[the nephron and filtration]
The nephron is the kidney's functional unit. Blood pressure drives fluid from the glomerular capillaries through a filtration membrane --- porous capillary endothelium, a basement membrane and the filtration slits of podocytes --- into the glomerular capsule. Water and small solutes pass; blood cells and most proteins do not. The tubule then reabsorbs and secretes substances using membrane transport proteins, which is where membrane structure connects to kidney function.
\end{background}

\begin{background}[renin and the juxtaglomerular apparatus]
The juxtaglomerular apparatus sits where the distal tubule meets the afferent arteriole. Its juxtaglomerular (granular) cells, in the arteriole wall, secrete renin when blood pressure or filtrate sodium falls, starting the renin--angiotensin--aldosterone system.
\end{background}

\begin{background}[oxytocin and labour]
Oxytocin is made in the hypothalamus and released from the posterior pituitary. During labour, stretching of the cervix triggers oxytocin release, which strengthens uterine contractions, which increase stretching: a positive feedback loop. Precise serum thresholds for triggering contractions are not the kind of figure introductory textbooks give, which is why the demo uses that question as an example the books cannot answer.
\end{background}

\begin{background}[chemotaxis]
Chemotaxis is the movement of a cell or organism towards or away from a chemical stimulus --- for example bacteria swimming towards nutrients, or white blood cells moving towards signals from damaged tissue.
\end{background}

\begin{background}[cardiac muscle and intercalated discs]
Cardiac muscle cells are joined end to end by intercalated discs, which contain desmosomes that hold the cells together under mechanical stress and gap junctions that let electrical signals pass directly between cells, so the heart contracts as a coordinated unit.
\end{background}

\begin{background}[crossing over]
During prophase I of meiosis, homologous chromosomes pair and non-sister chromatids exchange segments. The recombinant chromosomes carry new combinations of alleles, such as the non-parental genotypes in a genetics question.
\end{background}

\begin{background}[smooth muscle latch state]
Smooth muscle can hold tension for long periods with little energy use, in what is called the latch state, in which cross-bridges detach slowly. Some molluscs have a comparable \enquote{catch} mechanism. Comparing the two in detail goes beyond introductory textbooks, which is why the related gold-set question counts as unanswerable.
\end{background}
